\section{Gestión del proyecto}\label{sec:gestion}


\subsection{Aplicación de los principios del PMBOK}


La organización del equipo se apoya en los principios del PMBOK, séptima edición~\cite{pmbok2021}. Se asignan responsabilidades principales y se mantiene la revisión conjunta de los entregables. Esta distribución adapta el trabajo al tamaño del equipo y al contexto académico del proyecto; sus denominaciones no son cargos obligatorios del PMBOK.

\begin{itemize}
\item \textbf{Equipo:} los tres integrantes colaboran y conocen el contenido de los entregables, aunque cada uno coordina actividades diferentes.
\item \textbf{Interesados:} Jordan coordina las sesiones y la comunicación con el taller. Las interpretaciones que requieran confirmación se consultarán con los participantes pertinentes.
\item \textbf{Calidad:} Rai prepara los documentos y modelos; Jonathan revisa su claridad, coherencia y respaldo antes de cerrar cada versión.
\item \textbf{Adaptación al contexto:} el equipo ajustará las actividades a la disponibilidad del taller y a la información que pueda obtenerse, registrando los cambios acordados.
\item \textbf{Administración responsable:} Jordan administra la carpeta privada; todos respetan los consentimientos y evitan publicar información identificable.
\end{itemize}

Estas acciones constituyen una aplicación prevista de los principios, que deberá comprobarse mediante los aportes, revisiones, acuerdos y evidencias del proyecto.


\subsection{Roles y responsabilidades}


\begin{table}[H]
\centering\small
\caption{Roles y responsabilidades del equipo ContX.}\label{tab:roles}
\noindent\begin{tabularx}{\linewidth}{@{}L{3.6cm}L{3.2cm}Y@{}}
\toprule
\textbf{Integrante} & \textbf{Responsabilidad principal} & \textbf{Actividades} \\
\midrule
Holguin Mendoza Jonathan Alexander & Validación y trazabilidad & Revisar documentos y modelos; comprobar claridad, consistencia y verificabilidad; relacionar hallazgos, objetivos, alcance y evidencias; registrar observaciones y verificar las correcciones. \\
Alvarez Rubio Rai Eduardo & Análisis y documentación & Organizar y analizar la información obtenida; redactar los documentos; elaborar los modelos e incorporar las correcciones acordadas. \\
Basurto Anchaguano Jordan Steven & Coordinación y elicitación & Coordinar el contacto y las sesiones con el taller; preparar y conducir la elicitación con apoyo del equipo; gestionar aceptación y consentimientos; custodiar las evidencias privadas. \\
\bottomrule
\end{tabularx}
\end{table}


Las responsabilidades no son exclusivas: todos contribuyen y revisan el trabajo.

\subsection{Procedimiento de revisión y validación}
\label{sec:revision-entregables}

El responsable de validación y trazabilidad coordinará la revisión antes de cada entrega o actualización de un instrumento, documento o modelo:

\begin{enumerate}
\item \textbf{Delimitar la revisión.} Identificar el archivo y su versión, las fuentes disponibles y los productos que solicita la guía.
\item \textbf{Revisar el contenido.} Comprobar que las preguntas sean claras, neutrales y coherentes con el consentimiento; que visión, objetivos, alcance, partes interesadas, técnicas e instrumentos concuerden; y que los modelos y las conclusiones tengan respaldo. En los objetivos, revisar indicador, situación inicial, meta y plazo.
\item \textbf{Comprobar la trazabilidad.} Verificar las relaciones entre fuentes, hallazgos, objetivos, alcance y modelos según el criterio descrito a continuación.
\item \textbf{Comunicar observaciones.} Indicar el elemento afectado, el problema y la corrección propuesta. Separar los ajustes documentales de las dudas que deben consultarse al taller y revisar que no se publiquen datos personales.
\item \textbf{Verificar las correcciones.} El responsable del contenido realiza los ajustes y el de validación comprueba su resolución. Si cambian las necesidades, prioridades o el alcance, Jordan coordina su confirmación con el taller; la revisión interna no sustituye ese acuerdo.
\item \textbf{Cerrar la revisión.} Registrar lo revisado, las correcciones y los pendientes. Comprobar las citas bibliográficas, la compilación y las referencias a tablas y figuras antes de integrar la versión al repositorio.
\end{enumerate}

Las observaciones se marcarán como pendientes, en corrección o resueltas. Los datos y acuerdos sin confirmar conservarán su estado pendiente.

\subsection{Criterio de trazabilidad para E1}

Se usarán los identificadores de preguntas, instrumentos, partes interesadas, objetivos, alcance y evidencias del documento. Cada relación indicará su fuente, la referencia concreta y el estado de confirmación, mediante las tablas, figuras y referencias del informe. El manifiesto permitirá identificar los originales privados.

La relación será: \textbf{fuente y evidencia $\rightarrow$ hallazgo $\rightarrow$ objetivo o capacidad del alcance $\rightarrow$ elemento del modelo}, cuando corresponda. Un hallazgo también puede justificar una restricción o exclusión. Los requisitos detallados se vincularán cuando se elaboren en las siguientes entregas.

\subsection{Cronograma hasta la semana 17}
% Semana 4 = 16-22 nov 2026. Marque con \x las semanas de cada actividad.
% Ajuste las actividades a las entregas reales del PFC del silabo.
La Tabla~\ref{tab:cronograma} propone la distribución del trabajo hasta la semana 17.
Las semanas de las actividades posteriores a E1 son una planificación del equipo y se ajustarán
a las fechas oficiales de las siguientes entregas.

\newcommand{\x}{\cellcolor{uteqverde!60}}
\begin{table}[H]
  \centering
  \caption{Cronograma del PFC (semanas 4 a 17).}
  \label{tab:cronograma}
  \small
  \setlength{\tabcolsep}{3pt}
  \begin{tabular}{@{}L{5.2cm}*{14}{|c}|@{}}
    \toprule
    \textbf{Actividad} & 4 & 5 & 6 & 7 & 8 & 9 & 10 & 11 & 12 & 13 & 14 & 15 & 16 & 17 \\
    \midrule
    E1: visión, alcance, contexto y plan & \x &  &  &  &  &  &  &  &  &  &  &  &  &  \\
    Completar fuentes y mediciones &  & \x & \x &  &  &  &  &  &  &  &  &  &  &  \\
    Analizar y especificar requisitos &  &  & \x & \x & \x & \x &  &  &  &  &  &  &  &  \\
    Modelar y revisar requisitos &  &  &  &  &  & \x & \x & \x & \x &  &  &  &  &  \\
    Validar con la organización &  &  & \x &  &  &  & \x &  &  &  & \x & \x &  &  \\
    Corregir y consolidar el informe &  &  &  &  &  &  &  &  &  & \x & \x & \x & \x &  \\
    Integrar la entrega final del PFC &  &  &  &  &  &  &  &  &  &  &  &  &  & \x \\
    \bottomrule
  \end{tabular}
\end{table}

La entrevista SE-01 corresponde a la semana 4; las sesiones de las semanas siguientes
son una previsión y aún deben acordarse con el taller.

\subsection{Plan de comunicaci\'on con la organizaci\'on}
La Tabla~\ref{tab:comunicacion} define qu\'e se comunica, a qui\'en, por qu\'e medio y con qu\'e frecuencia.

\begin{table}[H]
  \centering
  \caption{Plan de comunicaci\'on.}
  \label{tab:comunicacion}
  \noindent\begin{tabularx}{\linewidth}{@{}YL{2.4cm}L{2.4cm}L{2.2cm}L{2.4cm}@{}}
    \toprule
    \textbf{Qu\'e se comunica} & \textbf{Destinatario} & \textbf{Medio} & \textbf{Frecuencia} & \textbf{Responsable} \\
    \midrule
    Avance y dudas por resolver & Responsable del taller (PI-01) & Resumen breve; canal por acordar & Al cerrar una versión o surgir una decisión & Jordan coordina; equipo prepara \\
    Agenda y condiciones de las sesiones & Roles convocados; PI-01 como contacto del taller & Canal por acordar, sin datos personales & Antes de cada actividad & Jordan \\
    Validación del problema, prioridades, objetivos y alcance propuesto & Responsable del taller (PI-01) & Por acordar con el taller & Antes de cerrar la propuesta y ante cambios relevantes & Jordan y Jonathan \\
    \bottomrule
  \end{tabularx}
\end{table}

Para esta consulta, Jonathan preparará un resumen de la propuesta y los puntos por confirmar.
Jordan coordinará la fecha y el medio con el responsable del taller. Se registrarán los acuerdos,
las correcciones y las decisiones pendientes, indicando la fecha y la versión revisada. La
actividad queda prevista; los resultados se incorporarán después de realizarla.

\subsection{Registro de riesgos}\label{sec:registro-riesgos}

La Tabla~\ref{tab:riesgos} recoge los riesgos que pueden afectar la elaboración y validación de los requisitos. Se basa en la entrevista SE-01 y en los datos que aún faltan por confirmar. Las respuestas son acciones previstas.

Se propone valorar la probabilidad (P) y el impacto (I) de 1 a 3: bajo, medio y alto. El producto P$\times$I servirá para ordenar la atención. Las valoraciones son estimaciones iniciales para revisión del equipo.

\begingroup
\small
\setstretch{1.1}
\setlength{\tabcolsep}{4pt}
\renewcommand{\arraystretch}{1.2}
\begin{longtable}{@{}L{0.9cm}L{4.5cm}cccL{4.7cm}L{2.5cm}@{}}
\caption{Riesgos del proyecto y respuestas previstas.}\label{tab:riesgos}\\
\toprule
\textbf{ID} & \textbf{Riesgo} & \textbf{P} & \textbf{I} & \textbf{P$\times$I} & \textbf{Respuesta} & \textbf{Responsable} \\
\midrule
\endfirsthead
\multicolumn{7}{@{}l}{\textit{Tabla~\ref{tab:riesgos}. Continuación.}}\\
\toprule
\textbf{ID} & \textbf{Riesgo} & \textbf{P} & \textbf{I} & \textbf{P$\times$I} & \textbf{Respuesta} & \textbf{Responsable} \\
\midrule
\endhead
\bottomrule
\endfoot
RI-01 & Formular objetivos que no puedan comprobarse, porque aún no hay cifras sobre la frecuencia de las dificultades ni el tiempo de búsqueda de información (C3 y D1). & 3 & 3 & 9 & Revisar los registros de ejemplo propuestos en E3 y acordar cómo obtener los datos que faltan. No completar cifras por suposición. Si no se obtienen, mantener la medición pendiente y explicar la limitación. & Análisis y documentación; revisión de validación y trazabilidad. \\
\addlinespace[7pt]
RI-02 & Retrasar las aclaraciones o la validación por depender del responsable del taller, quien también realiza toda la atención y reparación (A3 y D4). & 2 & 2 & 4 & Agrupar las dudas y coordinar con anticipación una conversación breve. Si no puede atenderla, continuar con los datos confirmados y dejar las decisiones pendientes para otra fecha acordada. & Coordinación y elicitación. \\
\addlinespace[7pt]
RI-03 & Convertir las dificultades mencionadas en funciones sin acordarlas con el taller, ampliando el alcance más allá de su prioridad de registro y control de equipos (C2 y C4 de SE-01). & 2 & 3 & 6 & Relacionar cada capacidad con una respuesta o evidencia y acordar con el responsable qué se incluirá. No presentar como confirmado lo que solo sea una propuesta del equipo. & Validación y trazabilidad; apoyo de análisis y documentación. \\
\addlinespace[7pt]
RI-04 & Publicar por descuido datos personales presentes en las notas, el consentimiento o los registros de ejemplo del taller, afectando la privacidad de participantes y clientes. & 2 & 3 & 6 & Mantener los originales en la carpeta privada y revisar la copia anonimizada antes de subirla. Si se detecta una publicación indebida, retirar el acceso público al archivo, avisar al docente y coordinar su tratamiento. & Custodia de evidencias; revisión de validación y trazabilidad. \\
\end{longtable}
\endgroup

El equipo revisará los riesgos durante el proyecto. Si alguno ocurre, el responsable aplicará la respuesta prevista.
